% GENERATED FILE. Edit canonical lessons, then run npm run generate:books.
% canonical-source-hash: fnv1a64:5925ce33fb85178a
% canonical-lessons: ML-C367-kurrakrtyam, ML-C367-kurrasammatam, ML-C367-pitikutuka, ML-C367-tattikkontupokuka, ML-C367-niriksikkuka

\chapter{A Crime and a Confession}
\label{ch:ml-a-crime-and-a-confession}
\hlchaptermodality{367}

\begin{chapteropening}
\textbf{By the end of this chapter:} \emph{I can say a crime, a confession, to capture, to kidnap and to keep watch on.}

Complete the last lesson of chapter 367: I can say a crime, a confession, to capture, to kidnap and to keep watch on.
\end{chapteropening}

\section[\texorpdfstring{\ml{കുറ്റകൃത്യം} (kuṟṟakṛtyaṁ)}{കുറ്റകൃത്യം (kuṟṟakṛtyaṁ)}]{\ml{കുറ്റകൃത്യം} (kuṟṟakṛtyaṁ) — a crime}

\label{lesson:ML-C367-kurrakrtyam}



\begin{quote}
Before the new one: say the Malayalam for a fingerprint, then the Malayalam for mysterious.
\end{quote}

\subsection*{You'll want to know: \ml{കുറ്റകൃത്യം}}
\textbf{\ml{കുറ്റകൃത്യം}} — \emph{kuṟṟakṛtyaṁ} — \textquotedblleft{}a crime\textquotedblright{}.

\begin{quote}
\ml{നഗരത്തിൽ} \ml{കുറ്റകൃത്യങ്ങൾ} \ml{കുറഞ്ഞു}.
\end{quote}

\begin{quote}
\emph{nagarattil kuṟṟakṛtyaṅṅaḷ kuṟaññu.}
\end{quote}

\begin{quote}
Crime fell in the city.
\end{quote}

Plural: \ml{കുറ്റകൃത്യങ്ങൾ}, \emph{kuṟṟakṛtyaṅṅaḷ}.

\ml{കുറ്റം}, an offence, + \ml{കൃത്യം}, a deed: a criminal act.

\begin{cousinweb}
\textbf{\ml{കുറ്റകൃത്യം}} (\emph{kuṟṟakṛtyaṁ}) joins the native \ml{കുറ്റം} (\emph{kuṟṟaṁ}), \textquotedblleft{}an offence\textquotedblright{} (DEDR 1851), and Sanskrit \emph{kṛtya}, \textquotedblleft{}to be done; a deed\textquotedblright{} (Monier-Williams).
\end{cousinweb}

\begin{grammarlens}[title={What you've built}]
One more word for telling what happened.
\end{grammarlens}

\subsection*{Your turn}
\begin{itemize}
\raggedright
  \item \emph{Say it:} \emph{kuṟṟakṛtyaṁ}
  \item \emph{Say it:} \emph{kuṟṟakṛtyaṁ}, once more
  \item \emph{Say it:} \emph{nagarattil kuṟṟakṛtyaṅṅaḷ kuṟaññu.} — Crime fell in the city.
  \item \emph{Recall:} say the Malayalam for a fingerprint, then the Malayalam for mysterious, then say \emph{kuṟṟakṛtyaṁ} again
\end{itemize}

\begin{tcolorbox}[breakable,colback=teal!4,colframe=teal!35!black,title={Before you move on}]
What does \ml{കുറ്റകൃത്യം} mean? (\textquotedblleft{}A crime\textquotedblright{}.) Say it once more.
\end{tcolorbox}
\section[\texorpdfstring{\ml{കുറ്റസമ്മതം} (kuṟṟasammataṁ)}{കുറ്റസമ്മതം (kuṟṟasammataṁ)}]{\ml{കുറ്റസമ്മതം} (kuṟṟasammataṁ) — a confession}

\label{lesson:ML-C367-kurrasammatam}



\begin{quote}
Before the new one: say the Malayalam for mysterious, then the Malayalam for a crime.
\end{quote}

\subsection*{You'll want to know: \ml{കുറ്റസമ്മതം}}
\textbf{\ml{കുറ്റസമ്മതം}} — \emph{kuṟṟasammataṁ} — \textquotedblleft{}a confession\textquotedblright{}.

\begin{quote}
\ml{ഒടുവിൽ} \ml{കള്ളൻ} \ml{കുറ്റസമ്മതം} \ml{നടത്തി}.
\end{quote}

\begin{quote}
\emph{oṭuvil kaḷḷan kuṟṟasammataṁ naṭatti.}
\end{quote}

\begin{quote}
At last the thief made a confession.
\end{quote}

\ml{കുറ്റം}, an offence, + \ml{സമ്മതം}, consent: agreeing that you did it.

\begin{cousinweb}
\textbf{\ml{കുറ്റസമ്മതം}} (\emph{kuṟṟasammataṁ}) joins \ml{കുറ്റം} (\emph{kuṟṟaṁ}), \textquotedblleft{}an offence\textquotedblright{} (DEDR 1851), and \ml{സമ്മതം} (\emph{sammataṁ}), \textquotedblleft{}consent\textquotedblright{}, borrowed from Sanskrit \emph{sammata}, \textquotedblleft{}agreed\textquotedblright{} (CDIAL 12973).
\end{cousinweb}

\begin{grammarlens}[title={What you've built}]
One more word for telling what happened.
\end{grammarlens}

\subsection*{Your turn}
\begin{itemize}
\raggedright
  \item \emph{Say it:} \emph{kuṟṟasammataṁ}
  \item \emph{Say it:} \emph{kuṟṟasammataṁ}, once more
  \item \emph{Say it:} \emph{oṭuvil kaḷḷan kuṟṟasammataṁ naṭatti.} — At last the thief made a confession.
  \item \emph{Recall:} say the Malayalam for mysterious, then the Malayalam for a crime, then say \emph{kuṟṟasammataṁ} again
\end{itemize}

\begin{tcolorbox}[breakable,colback=teal!4,colframe=teal!35!black,title={Before you move on}]
What does \ml{കുറ്റസമ്മതം} mean? (\textquotedblleft{}A confession\textquotedblright{}.) Say it once more.
\end{tcolorbox}
\section[\texorpdfstring{\ml{പിടികൂടുക} (piṭikūṭuka)}{പിടികൂടുക (piṭikūṭuka)}]{\ml{പിടികൂടുക} (piṭikūṭuka) — to capture; to catch}

\label{lesson:ML-C367-pitikutuka}



\begin{quote}
Before the new one: say the Malayalam for a crime, then the Malayalam for a confession.
\end{quote}

\subsection*{You'll want to know: \ml{പിടികൂടുക}}
\textbf{\ml{പിടികൂടുക}} — \emph{piṭikūṭuka} — \textquotedblleft{}to capture; to catch\textquotedblright{}.

\begin{quote}
\ml{നാട്ടുകാർ} \ml{കള്ളനെ} \ml{പിടികൂടി}.
\end{quote}

\begin{quote}
\emph{nāṭṭukār kaḷḷane piṭikūṭi.}
\end{quote}

\begin{quote}
The local people caught the thief.
\end{quote}

Past: \ml{പിടികൂടി}, \emph{piṭikūṭi}.

Common in the news for catching a thief or a wild animal.

\begin{cousinweb}
\textbf{\ml{പിടികൂടുക}} (\emph{piṭikūṭuka}) is a native Dravidian compound of \ml{പിടി} (\emph{piṭi}), \textquotedblleft{}a grip, a hold\textquotedblright{} (DEDR 4148), and \ml{കൂടുക} (\emph{kūṭuka}), \textquotedblleft{}to come together\textquotedblright{} (DEDR 1882): to close a grip on someone.
\end{cousinweb}

\begin{grammarlens}[title={What you've built}]
One more word for telling what happened.
\end{grammarlens}

\subsection*{Your turn}
\begin{itemize}
\raggedright
  \item \emph{Say it:} \emph{piṭikūṭuka}
  \item \emph{Say it:} \emph{piṭikūṭuka}, once more
  \item \emph{Say it:} \emph{nāṭṭukār kaḷḷane piṭikūṭi.} — The local people caught the thief.
  \item \emph{Recall:} say the Malayalam for a crime, then the Malayalam for a confession, then say \emph{piṭikūṭuka} again
\end{itemize}

\begin{tcolorbox}[breakable,colback=teal!4,colframe=teal!35!black,title={Before you move on}]
What does \ml{പിടികൂടുക} mean? (\textquotedblleft{}To capture; to catch\textquotedblright{}.) Say it once more.
\end{tcolorbox}
\section[\texorpdfstring{\ml{തട്ടിക്കൊണ്ടുപോകുക} (taṭṭikkoṇṭupōkuka)}{തട്ടിക്കൊണ്ടുപോകുക (taṭṭikkoṇṭupōkuka)}]{\ml{തട്ടിക്കൊണ്ടുപോകുക} (taṭṭikkoṇṭupōkuka) — to kidnap}

\label{lesson:ML-C367-tattikkontupokuka}



\begin{quote}
Before the new one: say the Malayalam for a confession, then the Malayalam for to capture.
\end{quote}

\subsection*{You'll want to know: \ml{തട്ടിക്കൊണ്ടുപോകുക}}
\textbf{\ml{തട്ടിക്കൊണ്ടുപോകുക}} — \emph{taṭṭikkoṇṭupōkuka} — \textquotedblleft{}to kidnap\textquotedblright{}.

\begin{quote}
\ml{രണ്ട്} \ml{പേർ} \ml{കച്ചവടക്കാരനെ} \ml{തട്ടിക്കൊണ്ടുപോയി}.
\end{quote}

\begin{quote}
\emph{raṇṭŭ pēr kaccavaṭakkārane taṭṭikkoṇṭupōyi.}
\end{quote}

\begin{quote}
Two men kidnapped the trader.
\end{quote}

Past: \ml{തട്ടിക്കൊണ്ടുപോയി}, \emph{taṭṭikkoṇṭupōyi}.

\ml{തട്ടി}, snatching, + \ml{കൊണ്ടുപോകുക}, to take away.

\begin{cousinweb}
\textbf{\ml{തട്ടിക്കൊണ്ടുപോകുക}} (\emph{taṭṭikkoṇṭupōkuka}) joins the native \ml{തട്ടുക} (\emph{taṭṭuka}), \textquotedblleft{}to strike, to knock\textquotedblright{} (DEDR 3039), and \ml{കൊണ്ടുപോകുക} (\emph{koṇṭupōkuka}), \textquotedblleft{}to take away\textquotedblright{}: to snatch someone and carry them off.
\end{cousinweb}

\begin{grammarlens}[title={What you've built}]
One more word for telling what happened.
\end{grammarlens}

\subsection*{Your turn}
\begin{itemize}
\raggedright
  \item \emph{Say it:} \emph{taṭṭikkoṇṭupōkuka}
  \item \emph{Say it:} \emph{taṭṭikkoṇṭupōkuka}, once more
  \item \emph{Say it:} \emph{raṇṭŭ pēr kaccavaṭakkārane taṭṭikkoṇṭupōyi.} — Two men kidnapped the trader.
  \item \emph{Recall:} say the Malayalam for a confession, then the Malayalam for to capture, then say \emph{taṭṭikkoṇṭupōkuka} again
\end{itemize}

\begin{tcolorbox}[breakable,colback=teal!4,colframe=teal!35!black,title={Before you move on}]
What does \ml{തട്ടിക്കൊണ്ടുപോകുക} mean? (\textquotedblleft{}To kidnap\textquotedblright{}.) Say it once more.
\end{tcolorbox}
\section[\texorpdfstring{\ml{നിരീക്ഷിക്കുക} (nirīkṣikkuka)}{നിരീക്ഷിക്കുക (nirīkṣikkuka)}]{\ml{നിരീക്ഷിക്കുക} (nirīkṣikkuka) — to keep watch on; to observe}

\label{lesson:ML-C367-niriksikkuka}



\begin{quote}
Before the new one: say the Malayalam for to capture, then the Malayalam for to kidnap.
\end{quote}

\subsection*{You'll want to know: \ml{നിരീക്ഷിക്കുക}}
\textbf{\ml{നിരീക്ഷിക്കുക}} — \emph{nirīkṣikkuka} — \textquotedblleft{}to keep watch on; to observe\textquotedblright{}.

\begin{quote}
\ml{പോലീസ്} \ml{ദിവസങ്ങളോളം} \ml{ആ} \ml{വീട്} \ml{നിരീക്ഷിച്ചു}.
\end{quote}

\begin{quote}
\emph{pōlīsŭ divasaṅṅaḷōḷaṁ ā vīṭŭ nirīkṣiccu.}
\end{quote}

\begin{quote}
The police watched that house for days.
\end{quote}

Past: \ml{നിരീക്ഷിച്ചു}, \emph{nirīkṣiccu}.

The noun is \ml{നിരീക്ഷണം}, observation.

\begin{cousinweb}
\textbf{\ml{നിരീക്ഷിക്കുക}} (\emph{nirīkṣikkuka}) is built from Sanskrit \emph{nirīkṣaṇa}, \textquotedblleft{}looking at, observing\textquotedblright{} (Monier-Williams), from \emph{nir-īkṣ}, \textquotedblleft{}to behold, to observe\textquotedblright{}, with the Malayalam verb ending -\ml{ിക്കുക}.
\end{cousinweb}

\begin{grammarlens}[title={What you've built}]
One more word for telling what happened.
\end{grammarlens}

\subsection*{Your turn}
\begin{itemize}
\raggedright
  \item \emph{Say it:} \emph{nirīkṣikkuka}
  \item \emph{Say it:} \emph{nirīkṣikkuka}, once more
  \item \emph{Say it:} \emph{pōlīsŭ divasaṅṅaḷōḷaṁ ā vīṭŭ nirīkṣiccu.} — The police watched that house for days.
  \item \emph{Recall:} say the Malayalam for to capture, then the Malayalam for to kidnap, then say \emph{nirīkṣikkuka} again
\end{itemize}

\begin{tcolorbox}[breakable,colback=teal!4,colframe=teal!35!black,title={Before you move on}]
What does \ml{നിരീക്ഷിക്കുക} mean? (\textquotedblleft{}To keep watch on; to observe\textquotedblright{}.) Say it once more.
\end{tcolorbox}
